\section{Discussion}
\label{sec:discussion}

\subsection{What the Doctrine Buys}

The confirmed clauses compose into a working control loop: verbalizable codes written through the band transport, at uncommitted moments, verified by band-targeted readback, inside a registered freedom budget. Each clause earns its place by a measured contrast. Removing coordinate discipline (writing borrowed geometry on a new plant) collapses lift from 0.40 to 0.05; removing timing discipline (rate-matched schedules) costs roughly half the gated advantage; removing the budget (continuous high-amplitude writing) quadruples the entropy cost; using the wrong verification instrument (a final-target lens) misreads loaded bands as empty. The doctrine is thus not a bundle of conventions but a set of independently falsifiable constraints, two of which failed in their first registered form and were revised under audit.

\subsection{Transparency Rather than Throughput}

The comparative evaluation forces a precise statement of what gated steering is for. If the goal is making a 4B model mention fire, continuous flooding achieves 0.96 surface rate at negligible entropy cost and requires no theory. The gated arm's 0.16 at the same cost would be a poor trade were throughput the objective. The objective the doctrine actually serves is different: placing chosen, verbalizable content into the model's workspace at moments when the model is receptive, verifying uptake through the model's own articulation, and bounding the intervention's effect on the model's continuation freedom. These properties, placement, verification, and budget, are what an auditor or a safety pipeline needs from a steering primitive, and none of them is provided by flooding. The finding also cautions against reading raw steering lift as an alignment-relevant capability measure.

\subsection{Loadability as a Model Property}

The loadability gradient (0.10 at 4B, 0.55 at 27B, 0.00 on the pruned twin) suggests that instructed workspace loading is an emergent capability with a size threshold inside a lineage, and that compression pipelines can remove it entirely while leaving band structure intact. The practical consequence runs in both directions: workspace steering of the kind studied here may become easier on larger models, and distillation may be a deliberate defence against it. Both directions are testable and neither is claimed beyond the three plants measured.

\subsection{The Research Loop as a Result}

The programme's methodological finding is that an expected-free-energy selector over registered menus, closed by dual gates and audited by isolated reviews, can run a multi-week interpretability programme with every decision replayable from committed artifacts. The audit trail did real work: the stream fix re-based every sampling estimate in the programme, the L18 supersession corrected stale dose levels by a measured $-0.10$, and a claimed replication was withdrawn when review showed it to be a deterministic replay artifact. These corrections are part of the released record, and the resulting ledger, loops L0 to L21 with twenty-seven selector cycles and eighteen reviews, is offered as a reusable template for small-compute rigorous interpretability work, complementing the intervention-level agent of the companion paper \cite{Sathish2026ACD}.

\subsection{Relation to the Source Philosophy}

The engineering reading extracted from Pratyabhij\~n\=a proved productive in a specific, limited sense: it generated a clause structure whose elements were independently falsifiable, two of which failed as first registered (the commitment-flash timing variant and any expectation of alignment dividends) and were revised or bounded accordingly. Nothing in the results bears on the philosophy's own claims about consciousness, and the paper asserts no such connection. The Sanskrit terminology is retained as precise internal nomenclature, in the same spirit as the companion works' usage \cite{Sathish2026ACD,Sathish2026Refusal}, with a complete engineering glossary in Appendix~\ref{app:glossary}.

\subsection{Limitations}
\label{sec:limitations}

Model scale and coverage are the primary limits: interventional results are established at 4B on two families, with 27B evidence restricted to loadability, so extrapolation of steering behaviour to larger models is unsupported. Behavioural metrics in the comparative evaluation are heuristic and pattern-based; judge-model scoring would strengthen the alignment-facing conclusions, and this limit is registered in the gates themselves. The contrastive experiments test single directions at a single site and single extraction seed, bounding naive application only; richer contrastive protocols could behave differently. Readback verification holds at screen tier on matched corpora and failed pooled generalisation, so the verification clause is the weakest confirmed link. Concept coverage is narrow (a handful of concrete nouns), corpus styles number three, and the corpus-amplitude interaction is supported at screen tier with one disclosed margin failure. The associative-store result is an integration proof under a cold store, and the trained-store comparison remains open. Finally, the compute envelope (18.17 GPU-hours on one shared node) bought multi-seed replication at small $n$ (15 to 40 per cell); several margins reported here, notably the alignment-advantage magnitude and the held-out readback balanced accuracy, sit inside confidence intervals that more compute would tighten.
